\documentclass[11pt]{article}
\usepackage[margin=1in]{geometry}
\usepackage{amsmath,amssymb}
\usepackage{times}

% Tighten spacing
\setlength{\parskip}{0.35em}
\setlength{\parindent}{0em}
\usepackage{titlesec}
\titlespacing*{\section}{0pt}{0.4em}{0.2em}

% Reduce space before title
\makeatletter
\def\@maketitle{
  \vspace*{-2em}
  \begin{center}
    {\Large\bfseries \@title \par}
    \vspace{0.5em}
    {\normalsize \@author \par}
    \vspace{0.25em}
    {\normalsize \@date \par}
  \end{center}
  \vspace{-1.2em}
}
\makeatother

\title{On the Intrinsic Dimensional Structure of Transformer Attention}
\author{Luis Chamberlain \\ \texttt{mcgrof.c@samsung.com}}
\date{December 5, 2025}

\begin{document}
\maketitle
\begin{center}
\textit{LayerNorm removes one dimension; learning determines the rest.}
\end{center}

\section*{Geometric Constraint}

LayerNorm centers each token vector by subtracting its mean component.  
This eliminates the $[1,1,\ldots,1]$ direction and places every input to attention on a $(d{-}1)$-dimensional hyperplane.  
All subsequent projections operate entirely within this reduced geometry.  
The geometry shows only that full rank is unnecessary; it does not predict how much lower the learned effective rank may become.

\section*{2D $\to$ 1D Example}

For $\mathbf{x} = [x_1,x_2]$ with mean $\mu = (x_1+x_2)/2$,
\[
\mathbf{z} = [x_1-\mu,\;x_2-\mu]
\]
satisfies
\[
z_1 + z_2 = 0 \quad\Rightarrow\quad z_2 = -z_1.
\]
Every centered vector lies on the line spanned by $[1,-1]$.  
Stacking $N$ such vectors produces a rank-1 matrix:
\[
\begin{bmatrix}
1.2 & -1.2\\
-0.8 & 0.8\\
2.0 & -2.0\\
0.5 & -0.5
\end{bmatrix}
=
\begin{bmatrix}1.2\\-0.8\\2.0\\0.5\end{bmatrix}
\begin{bmatrix}1\;\; -1\end{bmatrix}.
\]

\section*{Generalization to $d$ Dimensions}

LayerNorm enforces
\[
z_1 + z_2 + \cdots + z_d = 0,
\]
so all token embeddings lie in a $(d{-}1)$-dimensional subspace.  
Thus any matrix of centered embeddings satisfies
\[
\mathrm{rank}(Z) \le d{-}1.
\]
Since attention uses linear projections,
\[
Q = ZW_q,\qquad K = ZW_k,
\]
it follows that
\[
\mathrm{rank}(Q),\,\mathrm{rank}(K) \le d{-}1.
\]
All similarity structure in $QK^\top$ arises from these $(d{-}1)$ available directions.

\section*{Learned Dimensionality}

In practice, trained attention mechanisms often rely on a much smaller subset of directions inside the
LayerNorm hyperplane. The architectural constraint provides the ambient $(d{-}1)$-dimensional space; learning
determines how many of those dimensions are actually used. Thus the intrinsic geometry defines the ceiling,
while the effective rank emerges from optimization and data.

\end{document}
